\documentclass[10pt,a4paper]{amsart}
\usepackage[margin=19mm]{geometry}
\usepackage{amsmath,amssymb,mathtools}
\usepackage[scr=rsfs,cal=euler]{mathalfa}
\usepackage{xfrac}
\usepackage[unicode,hidelinks,bookmarks=false]{hyperref}
\newcommand{\ie}{i.e.\ }
\newcommand{\cf}{cf.\ }
\newcommand{\resp}{respectively\ }
\newcommand{\Subsection}{\subsection}
\providecommand{\nobreakdash}{\nobreak\hbox{-}}
\setlength{\emergencystretch}{2em}
\multlinegap0pt
\usepackage{amssymb}
\usepackage[shortlabels]{enumitem}
\usepackage{xcolor}
\newtheorem{proposition}{Proposition}
\usepackage{cleveref}
\crefname{proposition}{Proposition}{Propositions}
\newcommand{\lt}[1]{\texttt{\color{blue} TL: \footnotesize #1}}
\title{Deep learning and the rate of approximation by flows}
\author{Jingpu Cheng, Qianxiao Li, Ting Lin, Zuowei Shen}
\date{Source: arXiv:2603.15363v2 (CC BY 4.0)}
\begin{document}
\maketitle

\noindent\emph{Source and licence.} Deep learning and the rate of approximation by flows. Version arXiv:2603.15363v2 (CC BY 4.0), distributed by arXiv under the Creative Commons Attribution 4.0 International licence, http://creativecommons.org/licenses/by/4.0/, as recorded in the arXiv licence metadata for this exact version and shown on the abstract page https://arxiv.org/abs/2603.15363v2. Full licence notice: \texttt{licenses/arxiv-2603.15363-CC-BY-4.0.txt}.\par\smallskip

\noindent\textit{Editorial scope.} Counted unit: the article's proposition prop:interpolation\_lower\_bound (source0.tex lines 1256--1263). The article's complete original proof of it is delivered with it (source0.tex lines 1265--1440, one \texttt{proof} environment of 176 lines). No mathematics is written, repaired or re-derived: the statement and the proof are the article's own lines, in the article's own notation, and no retained line is substituted or re-flowed.  Retained uncounted as the source-local context the proof needs: lines 1001--1020.  Nothing was omitted from any retained span, so the package carries no omission record.  No retained line contains a citation, so the wrapper carries no bibliography and no bibliography entry.  No retained line contains a figure environment, an image-inclusion command or drawing code, so no image is delivered and the package declares no asset dependency.  Editorial formatting only, in addition: an AMS \texttt{amsart} preamble that restates the article's own declarations and macros used by the retained lines, namely the environments \texttt{proposition} on separate counters, together with the standard packages those lines need.  The retained environments print in this document as Proposition 1 in order of appearance, whereas the article prints the same items with the numbering of its own full document.  The complete native source of the article as distributed in the arXiv e-print for this exact version is bundled unchanged as \texttt{sources/196474/source0.tex}.
\begin{proposition}
\label{prop:min-Sn}
Let \begin{equation}
    \label{eq:k_i}
    \begin{aligned}
        k_i &= N\Delta^2 u(\frac{i}{N}):= N(u(\frac {i+1}{N})-2u(\frac i N)+u(\frac{i-1}{N})), \qquad i=1, \cdots, N-1,\\
    \end{aligned}
\end{equation}
and
\begin{equation}
    K_+ = \sum_{i=1}^{N-1} \max\{k_i, 0\}, \qquad K_- = \sum_{i=1}^{N-1} \min\{k_i, 0\}.
\end{equation}

Then, we have
\begin{equation}
\label{eq:min-Sn}
    \min S_N(w,v) = \sum_{i=1}^{N-1} |k_i| + \operatorname{dist}\Big(-N(u(\frac 1 N)-u(0)), [K_-, K_+] \Big).
\end{equation}

\end{proposition}

\begin{proposition}
    \label{prop:interpolation_lower_bound}
    For $ u \in \operatorname{BV}^2_0([0,1])$, the following identity holds
    \begin{equation}
        \inf \{ s \mid u \in \overline{\mathbf{CH}_s(\mathcal F)}^{C^0}\}
         \ge \|{u^\prime}\|_{\operatorname{TV}[0,1]}
    \end{equation}
\end{proposition}

\begin{proof}
% \lt{working/}
% \noindent\textbf{Step 1: optimal cost for $N$-grids interpolation}


% Since $h\in W^{2,1}([0,1])$, we have that


% Now we need to show that $\operatorname{Cost}(u)$ is also a lower bound for $\inf \{ s ~|~ u \in \overline{\mathbf{CH}_{s}(\mathcal F) \circ \psi}\}.$.

Suppose the opposite holds, then there exists a sequence
\begin{equation}
    g_k(x) = \sum_{i=1}^{N_k} \alpha_i\sigma(x-b_i)+\sum_{j=1}^{M_k} \beta_j\sigma(x-c_j), k=1,2, \cdots,
\end{equation}
such that
\begin{equation}
    \lim_{k\to \infty} g_k=u, \text{ and } \lim_{k\to \infty} \left(\sum_{i=1}^{N_k}|\alpha_i|+\sum_{j=1}^{M_k}|\beta_j|\right) = \|{u^\prime}\|_{\operatorname{TV}[0,1]}-\varepsilon,
\end{equation}
for some $\varepsilon >0$. We can assume that for each $k$, $\|g_k-u\|_{C([0,1])}\le \frac 1 k$ by choosing a proper subsequence, and $\left(\sum_{i=1}^{N_k}|\alpha_i|+\sum_{j=1}^{M_k}|\beta_j|\right) < \|{u^\prime}\|_{\operatorname{TV}[0,1]}.$

The key is to adjust the node points $b_i, c_i$ to some equidistributed points. To see this, we introduce the following modifications: For each $b_i\in [0,1]$, let $\tilde b_i$ be one of $\{0, \frac 1 k, \cdots, \frac{k-1}{k}, 1\}$ that is closest to $b_i$, for $i = 1,2,\cdots, N_k$.
Similarly, we define $\tilde c_i$.
For all $b_i<0$ and $c_j>1$, we rewrite $\alpha_i\sigma(x-b_i) = \alpha_i\sigma(x) - \alpha_i b_i$ and $\beta_j\sigma(c_j-x) = \beta_j\sigma(1-x) - \beta_j (c_i-1)$ for $x\in [0,1]$.
Subsequently, we define
\begin{equation}
    \begin{aligned}
            \tilde g_k(x) & := \sum_{i:b_i\in [0,1]} \alpha_i\sigma(x-\tilde b_i)+\sum_{j:c_j\in [0,1]} \beta_j\sigma(x-\tilde c_j)\\ &+\sum_{i:b_i<0} (\alpha_i\sigma(x) - \alpha_i b_i)  +\sum_{j:c_j>1} (\beta_j\sigma(1-x) - \beta_j (c_i-1))\\
            &= \sum_{i=0}^{k} \left(w_i \sigma(x-\frac{i}{k})+ v_i\sigma(\frac i k-x)\right) +C_k,
    \end{aligned}
\end{equation}
where $w_i$ and $v_i$ are the weights after merging, and $C_k$ is the constant term after merging.
% Notice that $\tilde g_k$ is linear in all the intervals $[\frac i k, \frac{i+1}{k}]$ for all $i=0,1 \cdots k-1$ and $\tilde g_k$ has a weight sum no more than that of $g_k$, i.e.
It then holds that
\begin{equation}
    \label{eq:weight_sum_upper_bound}
    \sum_{i=1}^k (|w_i|+|v_i|)\le \sum_{i=1}^{N_k}|\alpha_i|+\sum_{j=1}^{M_k}|\beta_j| \le \|{u^\prime}\|_{\operatorname{TV}[0,1]}-\varepsilon.
\end{equation}


Also, for sufficiently large $k$ we have:
\begin{equation}
    \|\tilde g_k-g_k\|_{C([0,1])}\le \frac{\|{u^\prime}\|_{\operatorname{TV}[0,1]}}{2} \cdot \frac{1}{k},
\end{equation}
thus
\begin{equation}
    \|\tilde g_k-u\|_{C([0,1])}\le \frac{1}{k}(\|{u^\prime}\|_{\operatorname{TV}[0,1]}+ 1).
\end{equation}


Now we fix $k$, and specify the error as
$\varepsilon_i = \tilde g_k(\frac i k)-u(\frac i k)$,
for $ i=0,1,\cdots, k,$ with  $|\varepsilon_j|\le  \|\tilde g_k-u\|_{C(K)}$.
Now we pivot at $\tilde g_k$, and the coefficients solve the following linear system:
\begin{equation}
    \label{eq:interpolation_perturbed}
    \frac{1}{N}(\sum_{i=j+1}^N iv_i+\sum_{i=0}^{j-1} (j-i)w_i)=u(\frac j N) -C_k+ \varepsilon_i, \qquad j=0,1, \cdots,N.
\end{equation}
In the following, we denote $u_i=u(\frac i k)$ for simplicity, and define:
\begin{equation}
    \begin{aligned}
            &\Delta^2 u_i: = u_{i+1}-2u_i+u_{i-1}, \quad i=1, \cdots, k-1, \quad \Delta^2 u_0: = u_1-u_0, \quad \Delta^2 u_k: = u_k-u_{k-1}.\\
            &\Delta^2 \varepsilon_i: = \varepsilon_{i+1}-2\varepsilon_i+\varepsilon_{i-1}, \quad i=1, \cdots, k-1, \quad\Delta^2 \varepsilon_0: = \varepsilon_1-\varepsilon_0, \quad \Delta^2 \varepsilon_k: = \varepsilon_k-\varepsilon_{k-1}.
    \end{aligned}
\end{equation}
Thanks to~\eqref{eq:interpolation_perturbed} and \Cref{prop:min-Sn}, we have:
\begin{equation}
    \label{eq:weight_sum_estimation_perturbed}
    \begin{aligned}
            \sum_{i=0}^{k} |w_i|+|v_i| \ge  k\sum_{i=1}^{k-1} |\Delta^2 u_i+\Delta^2 \varepsilon_i| + \operatorname{dist}\left( - k(\Delta^2 u_0+\Delta^2 \varepsilon_0),[\widetilde K_-, \widetilde K_+] \right)
    \end{aligned}
\end{equation}
where
\begin{equation}
    \widetilde K_+ = k \sum_{i=1}^{k-1} \max\{\Delta^2 u_i+\Delta^2 \varepsilon_i, 0\}, \qquad \widetilde K_- = k\sum_{i=1}^{k-1} \min\{\Delta^2 u_i+\Delta^2 \varepsilon_i, 0\}.
\end{equation}
The term $u_k$ and $\varepsilon_k$ do not appear explicitly in the terms of \eqref{eq:weight_sum_estimation_perturbed}. Nevertheless, we can perform a symmetrization for this discrete counterpart.
It follows from $k(\Delta^2 u_0+\Delta^2 \varepsilon_0)+\widetilde K_-+\widetilde K_{+} = k(\Delta^2 u_k+\Delta^2\varepsilon)$ that:
\begin{equation}
    \begin{aligned}
        &\operatorname{dist}\left( - k(\Delta^2 u_0+\Delta^2 \varepsilon_0),[\widetilde K_-, \widetilde K_+] \right) \\
        =&\operatorname{dist}\left( - k\left(\Delta^2 u_0+\Delta^2 \varepsilon_0\right)-\frac{\widetilde K_-+\widetilde K_+}{2},\left[\frac{\widetilde K_--\widetilde K_+}{2}, \frac{\widetilde K_+-\widetilde K_-}{2}\right] \right)\\
        = & \operatorname{dist}\left( -k\frac{\Delta^2 u_0+\Delta^2 \varepsilon_0+\Delta^2 u_k+\Delta^2\varepsilon_k}{2},\left[\frac{\widetilde K_--\widetilde K_+}{2}, \frac{\widetilde K_+-\widetilde K_-}{2}\right] \right)\\
        =& \frac{k}{2} \max\left\{|\Delta^2 u_0+\Delta^2 \varepsilon_0 + \Delta^2 u_k+\Delta^2 \varepsilon_k|-\sum_{i=1}^{k-1} |\Delta^2 u_i+\Delta^2 \varepsilon_i|, 0\right\}
    \end{aligned}
\end{equation}
Therefore, ~\eqref{eq:weight_sum_estimation_perturbed} can be rewritten in a symmetric form:
\begin{equation}
\label{eq:weight_sum_estimation_perturbed_symmetric}
 \sum_{i=0}^{k} (|w_i|+|v_i|)\ge \sum_{i=1}^{k-1} k|\Delta^2 u_i+\Delta^2 \varepsilon_i| + \frac{k}{2} \max\left\{|\Delta^2 u_0+\Delta^2 \varepsilon_0 + \Delta^2 u_k+\Delta^2 \varepsilon_k|-\sum_{i=1}^{k-1} |\Delta^2 u_i+\Delta^2 \varepsilon_i|, 0\right\},
\end{equation}

The rest of the proof is aimed at reducing the error term $\varepsilon_i$ in the right-hand side. In particular, we prove that for any $\delta > 0$
and sufficiently large $k$ we have

\begin{equation}
    \sum_{i=0}^k \left(|w_i|+|v_i|\right)\ge k\sum_{i=1}^{k-1} |\Delta^2 u_i|+ \frac 1 2 k\max \left\{ \left( |\Delta^2 u_0 + \Delta^2 u_k|-\sum_{i=1}^{k-1} |\Delta^2 u_i|\right), 0 \right\}- \delta.
\end{equation}


Now we start from \eqref{eq:weight_sum_estimation_perturbed_symmetric}. The first step is to relax $|x|$ and $\max\{x, 0\}$, namely,
\begin{equation}
    \label{eq:dual_representations}
    |x| = \max_{\phi\in [-1,1]} \phi x, \quad \max\{x, 0\} = \max_{\theta\in [0,1]} \theta x.
\end{equation}
We then have the following estimation:
\begin{equation}
    \begin{aligned}
        \text{RHS of }
        \eqref{eq:weight_sum_estimation_perturbed_symmetric}&\ge k\sum_{i=1}^{k-1} |\Delta^2 u_i+\Delta^2 \varepsilon_i|+ \frac \theta 2 k\left( |\Delta^2 u_0+\Delta^2 \varepsilon_0 + \Delta^2 u_k+\Delta^2 \varepsilon_k|-\sum_{i=1}^{k-1} |\Delta^2 u_i+\Delta^2 \varepsilon_i|\right)\\
        &=k\sum_{i=1}^{k-1} (1-\frac \theta 2) |\Delta^2 u_i+\Delta^2 \varepsilon_i| + \frac \theta 2 k|\Delta^2 u_0+\Delta^2 \varepsilon_0 + \Delta^2 u_k+\Delta^2 \varepsilon_k|\\
        & \ge k(1-\frac{\theta}{2})\sum_{i=1}^{k-1}  \phi_i(\Delta^2 u_i+\Delta^2 \varepsilon_i) + \frac \theta 2 k|\Delta^2 u_0+\Delta^2 \varepsilon_0 + \Delta^2 u_k+\Delta^2 \varepsilon_k|,
    \end{aligned}
\end{equation}
for all $\theta\in [0,1]$ and $\phi_i\in [-1,1]$ for $i=1, \cdots, k-1$.
For arbitrarily defined $\phi_0$ and $\phi_k$, the following discrete integration by parts hold:
\begin{equation}
    \sum_{i=1}^{k-1}\phi_i\Delta^2 \varepsilon_i =\sum_{i=1}^{k-1} \varepsilon_i \Delta^2\phi_i + \phi_{k-1} \varepsilon_k+\phi_1 \varepsilon_0-\varepsilon_{k-1}\phi_k-\varepsilon_1\phi_0,
\end{equation}
where $\Delta^2 \phi_i = \phi_{i+1}-2\phi_i+\phi_{i-1}$ is the discrete central difference. We then have the estimation:
\begin{equation}
    \begin{aligned}
            \sum_{i=1}^{k-1}  \phi_i(\Delta^2 u_i+\Delta^2 \varepsilon_i) &= \sum_{i=1}^{k-1} \phi_i\Delta^2 u_i+\sum_{i=1}^{k-1} \varepsilon_i \Delta^2\phi_i + \phi_{k-1} \varepsilon_k+\phi_1 \varepsilon_0-\varepsilon_{k-1}\phi_k-\varepsilon_1\phi_0\\
            &\ge \underbrace{\sum_{i=1}^{k-1} \phi_i\Delta^2 u_i- \frac{C}{k}\sum_{i=1}^{k-1} |\Delta^2\phi_i|}_{A} + \phi_{k-1} \varepsilon_k+\phi_1 \varepsilon_0-\varepsilon_{k-1}\phi_k-\varepsilon_1\phi_0.
    \end{aligned}
\end{equation}
Here, $C := \|{u^\prime}\|_{\operatorname{TV}[0,1]} + 1$ is independent with $k$.
Therefore, it follows that the right-hand side of \eqref{eq:weight_sum_estimation_perturbed_symmetric} is not less than
\begin{equation}
    \begin{aligned}
k(1 - \frac{\theta}{2}) A + \underbrace{k(1-\frac{\theta}{2}) \Big(\phi_{k-1} \varepsilon_k+\phi_1 \varepsilon_0-\varepsilon_{k-1}\phi_k-\varepsilon_1\phi_0 \Big) + \frac \theta 2 k|\Delta^2 u_0+\Delta^2 \varepsilon_0 + \Delta^2 u_k+\Delta^2 \varepsilon_k|.}_B
    \end{aligned}
\end{equation}


For a fixed $\delta > 0$, and $a = \operatorname{sign}(\Delta^2 \varepsilon_0+\Delta^2\varepsilon_1)$,  there exists a function $\phi \in C^2([0,1]),$ such that
\begin{enumerate}
        \item[(a)] $|\phi|_{C([0,1])}\le 1$;
        \item[(b)] $\phi(0) = \phi(1)=a$;
        \item[(c)] $\int_{[0,1]} \phi(x) u^{\prime\prime}(x)dx: = \int_{[0,1]}\phi d(Du^\prime) > \|u^\prime\|_{\operatorname{TV}([0,1])}-\delta.$
    \end{enumerate}

Fix $\delta$ and $a$, for each $k$, we choose $\phi_i = \phi(\frac i k)$ for $i=0,1,\cdots, k$. For the term $A$, we have
\begin{equation}
    \label{eq:main_term_estimate}
\begin{aligned}
    &\liminf_{k\to \infty}  k\sum_{i=1}^{k-1} \phi_i\Delta^2 u_i- C\sum_{i=1}^{k-1} |\Delta^2\phi_i| \\
    &=\liminf_{k\to \infty} \int_{[0,1]} \phi(x) u^{\prime\prime}(x)dx - \frac{C}{k}\int_{[0,1]} |\phi^{\prime\prime}(x)|dx\\
    & \ge \|u^\prime\|_{\operatorname{TV}([0,1])} -\delta.
\end{aligned}
\end{equation}
For the term $B$, we have
\begin{equation}
    \label{eq:remaining_term}
    \begin{aligned}
            &k(1-\frac \theta 2)(\phi_{k-1} \varepsilon_k+\phi_1 \varepsilon_0-\varepsilon_{k-1}\phi_k-\varepsilon_1\phi_0)+k\frac{\theta}{2}(|\Delta^2 u_0+\Delta^2 \varepsilon_0 + \Delta^2 u_k+\Delta^2 \varepsilon_k|)\\
            = & k(1-\frac \theta 2)\bigl(|\Delta^2 \varepsilon_0+\Delta^2\varepsilon_k|+ (a-\phi_{k-1})\varepsilon_k+(\phi_1-a)\varepsilon_0 \bigr)+k\frac{\theta}{2}(|\Delta^2 u_0+\Delta^2 \varepsilon_0 + \Delta^2 u_k+\Delta^2 \varepsilon_k|)\\
            \ge  & k(1-\frac \theta 2)|\Delta^2 \varepsilon_0+\Delta^2\varepsilon_k|+k\frac{\theta}{2}(|\Delta^2 u_0+ \Delta^2 u_k+\Delta^2 \varepsilon_0 +\Delta^2 \varepsilon_k|) - C(|\phi_{k-1}-a|+|\phi_1-a|)\\
            \ge & k\frac{\theta}{2}(|\Delta^2 \varepsilon_0+\Delta^2\varepsilon_k|+|\Delta^2 u_0+ \Delta^2 u_k+\Delta^2 \varepsilon_0 +\Delta^2 \varepsilon_k|)- C(|\phi_{k-1}-a|+|\phi_1-a|)\\
            \ge & k\frac \theta 2 |\Delta^2 u_0+\Delta^2 u_k|-C(|\phi_{k-1}-a|+|\phi_1-a|).
    \end{aligned}
\end{equation}
As $k\to \infty$, $\phi_{k-1}\to a$ and $\phi_1\to a$. Therefore, for sufficiently large $k$, we have that:
\begin{equation}
B \ge k\frac \theta 2 |\Delta^2 u_0+\Delta^2 u_k|-\delta.
\end{equation}
Combining all results above, we have shown that for sufficiently large $k$,
\begin{equation}
    \sum_{i=0}^k \left(|w_i|+|v_i|\right)\ge k\sum_{i=1}^{k-1} |\Delta^2 u_i|+ \frac \theta 2 k\left( |\Delta^2 u_0 + \Delta^2 u_k|-\sum_{i=1}^{k-1} |\Delta^2 u_i|\right)-\delta
\end{equation}
for all $\theta\in [0,1]$. Recall the dual form of $\max\{x, 0\}$ in~\eqref{eq:dual_representations},
we have shown that:
\begin{equation}
    \sum_{i=0}^k \left(|w_i|+|v_i|\right)\ge \|{u^\prime}\|_{\operatorname{TV}[0,1]} -3\delta.
\end{equation}
Since $\delta$ can arbitrarily small,  we then have a contradiction to~\eqref{eq:weight_sum_upper_bound}.
\end{proof}

\end{document}
